%! Least-Squares Regression — Unit 2, Lesson 2.4 — Slides (Beamer)
\documentclass[aspectratio=169, 11pt]{beamer}
\usepackage{apstats-beamer}   % provides \navyheader{} and \sectionlabel[color]{}

% The C and Y column types live in apstats-article, which a beamer deck does not
% load — redeclare them here rather than pulling in the article preamble.
\newcolumntype{C}[1]{>{\centering\arraybackslash}p{#1}}
\newcolumntype{Y}{>{\centering\arraybackslash}X}

\begin{document}

% TITLE SLIDE (CourseName is not defined in beamer, so the name is literal)
{
\setbeamercolor{background canvas}{bg=navy}
\begin{frame}[plain]
  \vspace{2.8cm}\hspace{0.55cm}%
  \begin{minipage}{0.9\textwidth}
    {\color{goldacc}\bfseries\small LESSON 2.4}\\[0.5cm]
    {\color{white}\bfseries\LARGE Least-Squares Regression}\\[1.2cm]
    {\color{skymid}\small Statistics~$\cdot$~Unit 2}
  \end{minipage}
\end{frame}
}

% GRADUAL RELEASE deck: title → targets → warm-up → I DO divider → one frame per notes
% instruction row (Least squares · Computer output & R-squared · Slope & intercept) → WE DO
% (Used cars) → spoken debrief → close & START THE HOMEWORK. No group activity, no exit ticket.
\newcommand{\redink}[1]{{\color{redacc}\bfseries #1}}

% ── LEARNING TARGETS ─────────────────────────────────────────────────────────
\begin{frame}[t, plain]
  \navyheader{Today's targets}
  \sectionlabel{Lesson 2.4}
  By the end of today, I can\ldots
  \begin{itemize}\setlength{\itemsep}{5pt}
    \item explain what makes the \textbf{least-squares regression line} the best line, and
          compute it from summary statistics.
    \item find the slope, intercept, and $r^2$ in \textbf{computer regression output}.
    \item interpret the \textbf{slope} and \textbf{$y$-intercept} in context.
    \item interpret the \textbf{coefficient of determination} $r^2$ in context.
  \end{itemize}
  \begin{block}{How today works}
    Warm-up $\to$ \textbf{notes together} $\to$ \textbf{one we work as a class} $\to$
    debrief $\to$ \textbf{you start the homework here, alone}.
  \end{block}
\end{frame}

% ── WARM-UP ──────────────────────────────────────────────────────────────────
\begin{frame}[t, plain]
  \navyheader{Warm-up}
  \sectionlabel{5 minutes $\cdot$ on your own, then we check}
  \begin{enumerate}\setlength{\itemsep}{6pt}
    \item $\widehat{\text{minutes}} = 8 + 1.5(\text{miles})$. Predict a 4-mile delivery. It
          took 16 minutes: what is its residual?
    \item Two lines, the same four points. Add each line's residuals. Then add their
          \emph{squares}. Which line fits better?
    \item $r = -0.80$. Describe it. Then compute $(-0.80)^2$ and $-0.80 \times 45/1.5$.
  \end{enumerate}
  \vspace{6pt}
  \begin{block}{Hold on to item 2}
    Both lines' residuals add to \redink{0}. So the plain sum cannot pick a winner ---
    the \emph{squares} can.
  \end{block}
\end{frame}

% ── I DO — direct instruction begins ─────────────────────────────────────────
{
\setbeamercolor{background canvas}{bg=navy}
\begin{frame}[plain]
  \vspace{1.8cm}\hspace{0.8cm}%
  \begin{minipage}{0.86\textwidth}
    {\color{goldacc}\bfseries\small GUIDED NOTES \ \textbullet\ \ I DO}\\[0.7cm]
    {\color{white}\bfseries\Large Open to your Guided Notes page. Fill each blank as we
    name it --- not before.}\\[0.9cm]
    {\color{skymid}\small Three terms go in the vocabulary box today. Write each one the
    moment we use it.}
  \end{minipage}
\end{frame}
}

% ── ROW 1: LEAST SQUARES ─────────────────────────────────────────────────────
\begin{frame}[t, plain]
  \navyheader{The least-squares regression line}
  \sectionlabel{notes, Least Squares $\cdot$ problems 1--3}
  \vspace{2pt}
  \begin{columns}[t]
    \begin{column}{0.48\textwidth}
      \begin{block}{The best line (EK DAT-1.G.1)}
        The one line that makes the sum of the \textbf{squared residuals} as small as
        possible. It always passes through $(\bar x, \bar y)$.
      \end{block}
      \vspace{4pt}
      {\small Twenty students: sleep $x$ (hr) and reaction time $y$ (ms).\\
      $\bar x = 7.2$, \ $s_x = 1.5$, \ $\bar y = 310$, \ $s_y = 45$, \ $r = -0.80$.}
    \end{column}
    \begin{column}{0.46\textwidth}
      \begin{enumerate}\setlength{\itemsep}{6pt}
        \item $b = r\,\dfrac{s_y}{s_x} = -0.80 \cdot \dfrac{45}{1.5} = -24$
        \item $a = \bar y - b\bar x = 310 - (-24)(7.2) = 482.8$
        \item $\widehat{\text{time}} = 482.8 - 24(\text{sleep})$
      \end{enumerate}
      \vspace{4pt}
      {\small $b$ always has the \redink{same sign} as $r$.}
    \end{column}
  \end{columns}
\end{frame}

% ── ROW 2: COMPUTER OUTPUT & R-SQUARED ───────────────────────────────────────
\begin{frame}[t, plain]
  \navyheader{Reading computer output}
  \sectionlabel{notes, Computer Output \& R-Squared $\cdot$ problems 4--6}
  \vspace{4pt}
  \begin{center}
  {\small\ttfamily
  \begin{tabular}{@{} l r r r r @{}}
  Predictor & Coef & SE Coef & T & P \\ \hline
  Constant & \redink{482.8} & 31.17 & 15.49 & 0.000 \\
  Sleep & \redink{$-$24.00} & 4.243 & $-$5.66 & 0.000 \\
  \end{tabular}\\[3pt]
  S = 27.74 \quad \redink{R-Sq = 64.0\%} \quad R-Sq(adj) = 62.0\%}
  \end{center}
  \vspace{2pt}
  \begin{columns}[t]
    \begin{column}{0.47\textwidth}
      \begin{block}{Where the line lives}
        $a$ = Coef on \textbf{Constant}. \ $b$ = Coef on the $x$ row (\textbf{Sleep}).
      \end{block}
    \end{column}
    \begin{column}{0.47\textwidth}
      \begin{block}{$r^2$ (EK DAT-1.G.4)}
        About 64\% of the \textbf{variation in reaction time} is explained by the linear model
        with sleep.
      \end{block}
    \end{column}
  \end{columns}
  \vspace{6pt}
  \begin{center}
    {\small $r = \pm\sqrt{0.64}$ --- the square root forgets the sign. The slope is negative,
    so $r = $ \redink{$-0.80$}.}
  \end{center}
\end{frame}

% ── ROW 3: SLOPE & INTERCEPT — THE CRUX ──────────────────────────────────────
\begin{frame}[t, plain]
  \navyheader{Saying what the numbers mean}
  \sectionlabel{notes, Slope \& Intercept $\cdot$ problems 7--10 $\cdot$ the point of today}
  \vspace{4pt}
  \begin{block}{Slope (EK DAT-1.H.2)}
    For each additional hour of sleep, the \redink{predicted} reaction time decreases by 24 ms.
  \end{block}
  \vspace{4pt}
  \begin{block}{Intercept (EK DAT-1.G.3, H.3)}
    Predicted 482.8 ms at 0 hours of sleep --- \textbf{not meaningful}: nobody in the data slept
    less than 4 hours.
  \end{block}
  \vspace{8pt}
  \begin{center}
    Maya: \emph{``When a student sleeps one more hour, their reaction time goes down 24 ms.''}\\[4pt]
    {\small\color{slate} Real students scatter around the line, and nobody assigned the sleep.
    The line promises a \textbf{prediction}, not an outcome --- and not a cause.}
  \end{center}
\end{frame}

% ── WE DO ────────────────────────────────────────────────────────────────────
\begin{frame}[t, plain]
  \navyheader{Guided Practice: Used Cars}
  \sectionlabel[goldacc]{We do $\cdot$ you hold the pen $\cdot$ problems 11--14}
  Thirty used Honda Civics, ages 0 to 10 years; price in thousands of dollars.
  \begin{center}
  {\small\ttfamily
  \begin{tabular}{@{} l r r r r @{}}
  Predictor & Coef & SE Coef & T & P \\ \hline
  Constant & 24.60 & 0.950 & 25.89 & 0.000 \\
  Age & $-$1.850 & 0.169 & $-$10.93 & 0.000 \\
  \end{tabular}\\[3pt]
  S = 2.10 \quad R-Sq = 81.0\% \quad R-Sq(adj) = 80.3\%}
  \end{center}
  \begin{itemize}\setlength{\itemsep}{3pt}
    \item Write the line. Interpret the slope --- say \emph{predicted}.
    \item Interpret the intercept. Is a 0-year-old Civic a real car?
    \item Find $r$ (which sign?) and interpret $r^2$.
  \end{itemize}
\end{frame}

% ── DEBRIEF ──────────────────────────────────────────────────────────────────
\begin{frame}[t, plain]
  \navyheader{Debrief: what did Maya get wrong?}
  \sectionlabel{8 minutes $\cdot$ pens down, papers up --- we say these aloud}
  \vspace{6pt}
  \begin{columns}[t]
    \begin{column}{0.46\textwidth}
      \begin{block}{Almost right}
        ``Each extra hour, reaction time goes down 24 ms.'' The number is right. The word
        that is missing is \redink{predicted}.
      \end{block}
    \end{column}
    \begin{column}{0.46\textwidth}
      \begin{block}{Two intercepts}
        Sleep: 0 hours is outside the data --- no meaning. Civics: age 0 is a new car ---
        \textbf{it means something}.
      \end{block}
    \end{column}
  \end{columns}
  \vspace{12pt}
  \begin{block}{Say it out loud}
    $r^2$ is about \redink{variation in $y$}, not about how many points are on the line, and
    not about ``percent correct.''
  \end{block}
\end{frame}

% ── CLOSE & START THE HOMEWORK (the individual practice) ─────────────────────
\begin{frame}[t, plain]
  \navyheader{Homework --- start it now}
  \sectionlabel[goldacc]{5 items $\cdot$ scored $\cdot$ due the first class after two study halls}
  Car weight and highway mileage --- a new context, on your own.
  \begin{itemize}\setlength{\itemsep}{4pt}
    \item \textbf{A1--A3:} build the line from summary statistics; interpret the slope and
          the intercept.
    \item \textbf{B1--B2:} interpret $r^2$; one multiple-choice item with a justification.
  \end{itemize}
  \begin{block}{Silent and alone --- I am circulating}
    Start with A1 and A2 while I can still catch a wrong start. The \textbf{AP Practice} page
    (iced drinks) is \redink{not scored} --- work it for the reps and bring it to review.
  \end{block}
\end{frame}

% ── CLOSE ────────────────────────────────────────────────────────────────────
{
\setbeamercolor{background canvas}{bg=navy}
\begin{frame}[plain]
  \vspace{1.5cm}\hspace{0.8cm}%
  \begin{minipage}{0.86\textwidth}
    {\color{goldacc}\bfseries\small BEFORE YOU GO}\\[0.7cm]
    {\color{white}\bfseries\Large The slope tells you how the prediction changes --- not what
    happens to every individual.}\\[0.9cm]
    {\color{skymid}\small Next: Lesson 2.5 --- departures from linearity. What one unusual
    point can do to $b$, $a$, and $r^2$.}
  \end{minipage}
\end{frame}
}

\end{document}
